
\newcommand{\city}{Санкт-Петербург} 
\newcommand{\university}{Национальный исследовательский университет ИТМО}  % первая строка
\newcommand{\department}{Факультет программной инженерии и компьютерной техники}  % Вторая строка
\newcommand{\major}{Направление программная инженерия}

\newcommand{\education}{Образовательная программа системное и прикладное программное обеспечение}  % четвертая строка

%<<<<< Информация о кафедре

%>>>>> Назание работы
\newcommand{\reporttype}{ОТЧЕТ ПО ПРОЕКТНОЙ РАБОТЕ} % тип работы, (главный заголовок титульного листа)
\newcommand{\lab}{Введение}        
\newcommand{\subject}{Физика}  
\newcommand{\labtheme}{Исследование влияния электромагнитных помех на сетевое оборудование в контролируемой среде} 
\newcommand{\student}{Шубин Егор Вячеславович\\Мухамедьяров Артур Альбертович}    
\newcommand{\studygroup}{P3209}                 % определение учебной группы 
\newcommand{\teacher}{
    Сорокина Елена Константиновна\\
}
%<<<<<<<<<<<<<<<<<<<<<<<<<< ПЕРЕМЕННЫЕ <<<<<<<<<<<<<<<<<<<<<<<<<<<<<<<<<<<


%>>>>>>>>>>>>>>>>>>>>>> ПРЕАМБУЛА >>>>>>>>>>>>>>>>>>>>>>>>>
%>>>>>>>>>>>>>>>>>> ПРЕАМБУЛА >>>>>>>>>>>>>>>>>>>>
\documentclass[14pt,final,oneside]{extreport}% класс документа, характеристики
%>>>>> Разметка документа
\usepackage[a4paper, mag=1000, left=3cm, right=1.5cm, top=2cm, bottom=2cm, headsep=0.7cm, footskip=1cm]{geometry} % По ГОСТу: left>=3cm, right=1cm, top=2cm, bottom=2cm,
\linespread{1} % межстройчный интервал по ГОСТу := 1.5
%<<<<< Разметка документа

%>>>>> babel c языковым пакетом НЕ должны быть первым импортируемым пакетом
\usepackage[utf8]{inputenc}
\usepackage[T1,T2A]{fontenc}
\usepackage[russian]{babel}
%<<<<<
\usepackage{amsmath, amssymb}
\usepackage{siunitx}
\usepackage{enumitem}
\usepackage{geometry}
\usepackage{caption}

%>>>...>> прочие полезные пакеты
\usepackage{amsmath,amsthm,amssymb}
\usepackage{mathtext}
\usepackage{braket}
\usepackage{indentfirst}
\usepackage{graphicx}
\usepackage{float}
\graphicspath{{/home/ivan/itmo/informatics/latex}}
\DeclareGraphicsExtensions{.pdf,.png,.jpg}
%\usepackage{bookmark}
\usepackage{courier}
\usepackage[dvipsnames]{xcolor}
\usepackage{hyperref}  % Использование ссылок
\hypersetup{%  % Настройка разметки ссылок
    colorlinks=true,
    linkcolor=blue,
    filecolor=magenta,      
    urlcolor=magenta,
    %pdftitle={Overleaf Example},
    %pdfpagemode=FullScreen,
}
\usepackage{pgfplots}

\usepackage{multicol} 
\usepackage[table]{xcolor}

\setlength{\columnsep}{.15\textwidth} 
\usepackage{listings} 
\usepackage{caption}
\DeclareCaptionFont{white}{\color{white}} 
\DeclareCaptionFormat{listing}{\colorbox{gray}{\parbox{\textwidth}{#1#2#3}}}


\makeatletter
\newcommand{\iffby}[2][]{\ext@arrow 0359\Rightarrowfill@{#1}{#2}}% равносильность с надписью
\makeatother




\newcommand\Section[1]{
    % Принимает 1 аргумент - название секции
    \refstepcounter{section}
    \section*{%
        \raggedright
        % Отключена дополнительная нумерация chapter в section в тексте документа:
        % \arabic{chapter}.\arabic{section}. #1}
        % Отключена любая нумарация section в тексте документа: 
        \arabic{section}. #1%
    }
    
    % Отключена дополнительная нумерация chapter в section в toc (table of contents) Оглавлении (Содержании):
    % \addcontentsline{toc}{section}{\arabic{chapter}.\arabic{section}. #1}
    \addcontentsline{toc}{section}{\arabic{section}. #1} 
}

\newcommand\Chapter[3]{%
    % Принимает 3 аргумента - название главы и дополнительный заголовок и множитель ширины загловка (можно ничего)
    \refstepcounter{chapter}%
    \chapter*{%
        %\hfill % заполнение отступом пространства до заголовка
        \begin{minipage}{#3\textwidth} % Можно изменить ширину министраницы (заголовка)
            \flushleft % Выранивание заголовка по левому краю параграфа (заголовка)
            %\flushright % Выранивание заголовка по правому краю параграфа (заголовка)
            \begin{huge}%
                % Отключена нумерация глав в тексте:
                % \textbf{\chaptername\ \arabic{chapter}\\}
                \textbf{#1}% Первый заголовок
            \end{huge}%
            \\% Перенос сторки
            \begin{Huge}
                #2% Второй заголовок
            \end{Huge}
        \end{minipage}
    }%
    % Отключена нумерация для chapter в toc (table of contents), т.е. Оглавлении (Содержании):
    % \addcontentsline{toc}{chapter}{\arabic{chapter}. #1}
    % Представление главы в содержании:
    \addcontentsline{toc}{chapter}{#1. #2}%
}
\newcommand\Subsection[1]{
    % Принимает 1 аргумент - название подсекции
    \refstepcounter{subsection}
    \subsection*{%
        \raggedright%
        % Отключена дополнительная нумерация chapter в section в тексте документа (можно добавить отступ с помощью :
        % \arabic{chapter}.\arabic{section}.\arabic{subsection}. #1}
        \hspace*{12pt}
        \arabic{section}. \arabic{subsection}. #1
    }
    % Отключена дополнительная нумерация chapter в section в Оглавлении (Содержании):
    %\addcontentsline{toc}{subsection}{\arabic{chapter}.\arabic{section}.\arabic{subsection}. #1}
    \addcontentsline{toc}{subsection}{\arabic{subsection}. #1}
}
\newcommand\Subsubsection[1]{
    % Принимает 1 аргумент - название подсекции
    \refstepcounter{subsubsection}
    \subsubsection*{%
        \raggedright%
        % Отключена дополнительная нумерация chapter в section в тексте документа (можно добавить отступ с помощью :
        % \arabic{chapter}.\arabic{section}.\arabic{subsection}. #1}
        \hspace*{12pt}
        \arabic{section}. \arabic{subsection}. \arabic{subsubsection}. #1
    }
    % Отключена дополнительная нумерация chapter в section в Оглавлении (Содержании):
    %\addcontentsline{toc}{subsection}{\arabic{chapter}.\arabic{section}.\arabic{subsection}. #1}
    \addcontentsline{toc}{subsubsection}{\arabic{subsubsection}. #1}
}


\lstdefinestyle{customasm}{
    belowcaptionskip=1\baselineskip,
    frame=single, 
    frameround=tttt,
    xleftmargin=\parindent,
    language=[x86masm]Assembler,
    basicstyle=\footnotesize\ttfamily,
    commentstyle=\itshape\color{green!60!black},
    keywordstyle=\color{blue!80!black},
    identifierstyle=\color{red!80!black},
    tabsize=4,
    numbers=left,
    numbersep=8pt,
    stepnumber=1,
    numberstyle=\tiny\color{gray}, 
    columns = fullflexible,
    morekeywords={ORG, LD, ST, POP, SWAB, JUMP, BEQ, SXTB, CLA}
}










%<<<<<<<<<<<<<<<<<<<<<< ПРЕАМБУЛА <<<<<<<<<<<<<<<<<<<<<<<<<


%%%%%%%%%%%%%%%%%%% СОДЕРЖИМОЕ ОТЧЕТА %%%%%%%%%%%%%%%%%%%%%
%>>>>>>>>>>>>>>> ''''''''''''''''''''''' >>>>>>>>>>>>>>>>>>
\begin{document}

\renewcommand{\chaptername}{\lab\ \labnumber} % переименование глав
\def\contentsname{Содержание} 
\begin{titlepage}

  % Название университета
  \begin{center}
    \textsc{%
      \university\\[5mm]
      \department\\[2mm]
      \major\\
      \education\\
      \specialization\\
    }

    \vfill
    % Название работы
    \textbf{\reporttype\\[3mm]
      курса <<\subject>> \\[6mm]
      по теме: <<\labtheme>>\\[3mm]
    }
  \end{center}

  \vfill
  \hfill
  % Информация об авторе работы и проверяющем
  \begin{minipage}{.5\textwidth}
    \begin{flushright}

      Выполнили студенты:\\[2mm]
      \student\\[2mm]
      группа: \studygroup\\[5mm]

      Преподаватель:\\[2mm]
      \teacher

    \end{flushright}
  \end{minipage}

  \vfill

  % Нижний колонтитул первой страницы
  \begin{center}
    \city, \the\year\,г.
  \end{center}

\end{titlepage}
%<<<<<<<<<<<<<<<<<<< ТИТУЛЬНЫЙ ЛИСТ <<<<<<<<<<<<<<<<<<<<<<<


\tableofcontents

\setcounter{page}{2}
\newpage
\Section{Проектная работа}

\Subsection{Аннотация}

Целью данного проекта является разработка прототипа автоматической
системы тестирования устойчивости сетевого оборудования к
электромагнитным помехам (EMI).

Современные сетевые устройства работают в условиях высокой плотности
электронных систем и источников электромагнитного излучения.
Электромагнитные помехи могут негативно влиять на работу оборудования,
вызывая потерю пакетов, увеличение задержек, нестабильность соединения
или снижение пропускной способности сети.

В рамках проекта разрабатывается устройство на базе микроконтроллера
ESP32, генерирующее контролируемые электромагнитные помехи заданной
частоты и мощности посредством схемы
MOSFET~+~индуктивная катушка. Параллельно с генерацией помех
внешний EMI-датчик непрерывно измеряет реальный уровень
электромагнитного поля в контролируемой зоне,
фиксируя не только излучение генератора,
но и фоновые помехи от окружающих устройств
(блоки питания, зарядные устройства, кабели под нагрузкой).

Управляющий Python-скрипт автоматически изменяет параметры генератора,
одновременно запуская сетевые тесты с помощью инструментов
\texttt{iperf3}, \texttt{ping} и \texttt{speedtest}.
Данные с EMI-датчика и результаты сетевых тестов
синхронизируются по временной метке и анализируются в реальном времени.
По результатам измерений строится трёхмерный график,
отражающий корреляцию между реальным уровнем электромагнитного поля
и деградацией сетевых параметров: потерей пакетов, джиттером
и пропускной способностью.

Система позволяет автоматически определять порог устойчивости
сетевого оборудования к электромагнитным воздействиям
и может применяться для сравнительного тестирования кабелей
(UTP, STP), роутеров и другого сетевого оборудования
в условиях контролируемых электромагнитных помех.


\Subsection{Цель проекта}

Разработать прототип автоматической системы,
которая генерирует управляемые электромагнитные помехи
и одновременно измеряет их реальный уровень с помощью EMI-датчика,
проводя сетевые тесты в режиме реального времени
и выявляя количественную зависимость между уровнем электромагнитного
воздействия и параметрами работы сетевого оборудования.

\Subsection{Задачи проекта}

\begin{enumerate}

\item Изучить физические принципы возникновения
электромагнитных помех, механизмы их влияния на
электронные и сетевые устройства, а также существующие
методы измерения электромагнитного поля.

\item Разработать аппаратную часть генератора помех
на базе микроконтроллера ESP32 со схемой MOSFET~+~катушка,
обеспечивающей изменение частоты ШИМ-сигнала
в рабочем диапазоне с регулируемым уровнем мощности

\item Интегрировать внешний EMI-датчик для непрерывного
измерения реального уровня электромагнитного поля
в контролируемой зоне, включая фоновые помехи
от сторонних устройств.

\item Написать управляющий Python-скрипт,
который автоматически изменяет параметры генератора
и запускает сетевые тесты (\texttt{iperf3}, \texttt{ping},
\texttt{speedtest}), синхронизируя результаты с показаниями
EMI-датчика по временной метке.

\item Реализовать систему визуализации данных:
построение трёхмерного графика зависимости
реального уровня EMI от параметров сети
(потеря пакетов, джиттер, пропускная способность).

\item Провести сравнительное экспериментальное исследование
устойчивости различных типов сетевых кабелей и оборудования
(например UTP или STP, роутер, PoE-устройства) к контролируемым
электромагнитным воздействиям и определить пороговые
значения помех, при которых начинается деградация соединения.

\end{enumerate}
\end{document}
